% ID: SEL-D0017-BA-E02
% ESTADO: cotejada. Segunda pasada visual completa y borrador compilado revisados.
% Fuente: Selectividad/matematicas_aplicadas/matematicas_aplicadas_P_E.pdf
% SHA-256: bc88bfeb0f184eb0c806c4efe04e84c24f3e5fcef9caaa12e94fd48d9f67a389
% Página: 1; bloque A; ejercicio original 2.

Se consideran las matrices
$A=\begin{pmatrix}1&0&1\\k&-3&2\\1&k&1\end{pmatrix}$,
$B=\begin{pmatrix}3&-2&-1\\0&1&1\end{pmatrix}$ y
$C=\begin{pmatrix}-2&4\\1&3\end{pmatrix}$.
\begin{enumerate}[label=\alph*)]
  \item \textbf{(0.8 puntos)} Razone si las siguientes operaciones se pueden
  realizar y en aquellos casos en que sea posible, indique la dimensión
  de la matriz resultante:
  \[
  B^t\cdot A\qquad C\cdot B\qquad B\cdot A+B\qquad B^2
  \]
  \item \textbf{(0.7 puntos)} Calcule los valores del parámetro $k$ para
  los que la matriz $A$ es invertible.
  \item \textbf{(1 punto)} Para $k=-1$, calcule la inversa de la matriz $A$.
\end{enumerate}
